%=============================================================================
% COMPLETE DERIVATION CHAIN: FIRST PRINCIPLES --> alpha^{-1} = 137.03599985
% Extracted equation-by-equation from section_alpha_locked.tex and appendix_K.tex
% With all gaps identified
% Generated 2026-03-22
%=============================================================================
\documentclass[12pt,a4paper]{article}
\usepackage{amsmath,amssymb,amsthm}
\usepackage{booktabs}
\usepackage{enumitem}
\usepackage[dvipsnames]{xcolor}
\usepackage{tcolorbox}
\usepackage{geometry}
\geometry{margin=2.5cm}
\usepackage{hyperref}

\newtheorem{step}{Step}
\newtheorem{gap}{GAP}
\newcounter{eqtrace}
\newcommand{\eqstep}[1]{\refstepcounter{eqtrace}\noindent\textbf{Equation \arabic{eqtrace}:} #1\par\medskip}

\title{Complete Equation-by-Equation Derivation Chain\\[6pt]
\large From First Principles to $\alpha^{-1} = 137.03599985$\\[4pt]
\normalsize Extracted from \texttt{section\_alpha\_locked.tex}
and \texttt{appendix\_K.tex}}
\author{Trace Audit --- Density Field Dynamics v108}
\date{22 March 2026}

\begin{document}
\maketitle

\tableofcontents
\newpage

%=============================================================================
\section{Overview: The Two Routes}
%=============================================================================

The DFD derivation of $\alpha$ operates through \textbf{two routes}
of different precision, both anchored to the same master formula:

\begin{equation}\tag{MASTER}
\boxed{\alpha^{-1} = 6\pi\,\kappa_{U(1)}}
\end{equation}

\begin{itemize}
\item \textbf{Route 1 (Chern--Simons + Lattice):} Predicts the bare lattice
  coupling $\beta_{U(1)} = 3.7969$ from topological inputs, runs lattice
  Monte Carlo to measure $\kappa_{U(1)}$, obtains $\alpha^{-1} \approx 137$
  at $\sim 1\%$ precision.

\item \textbf{Route 2 (Convention-Locked Spectral Action):} Computes
  $\kappa_{U(1)} = 7.26999\ldots$ analytically from the $a_4$ Seeley--DeWitt
  coefficient on Toeplitz-truncated $\mathbb{C}P^2 \times S^3$, obtaining
  $\alpha^{-1} = 137.03599985$ at sub-ppm precision.
\end{itemize}

Both routes share Steps 1--11 below.  They diverge at Step 12
(lattice measurement vs.\ spectral action computation).

%=============================================================================
\section{Route 1: Topology $\to$ Chern--Simons $\to$ Lattice $\to$ $\alpha$}
\label{sec:route1}
%=============================================================================

%---------------------------------------------------------------------
\subsection{Step 1: Internal manifold}
%---------------------------------------------------------------------

\textbf{Postulate:} The internal (microsector) geometry is the 7-manifold
\begin{equation}\tag{E1}
\mathcal{M}_7 = \mathbb{C}P^2 \times S^3.
\end{equation}

\textit{Source:} Core DFD gauge-emergence framework (Appendix F).
This is the starting geometric input; everything downstream follows from
the topology and Riemannian structure of this space.

%---------------------------------------------------------------------
\subsection{Step 2: Canonical bundle of $\mathbb{C}P^2$}
%---------------------------------------------------------------------

\textbf{Theorem (algebraic geometry):} The canonical line bundle of
$\mathbb{C}P^2$ is:
\begin{equation}\tag{E2}
K_{\mathbb{C}P^2} = \mathcal{O}(-3).
\end{equation}

\textit{Status:} Rigorous mathematical theorem.  This is the first Chern
class of the anti-canonical bundle; $c_1(\mathbb{C}P^2) = 3$.

%---------------------------------------------------------------------
\subsection{Step 3: Spin$^c$ structure and determinant line bundle}
%---------------------------------------------------------------------

The canonical Spin$^c$ structure on $\mathbb{C}P^2$ has determinant
line bundle:
\begin{equation}\tag{E3}
L_{\det} = K^{-1}_{\mathbb{C}P^2} = \mathcal{O}(3).
\end{equation}

\textit{Status:} Rigorous.  $\mathbb{C}P^2$ is not spin but admits a
canonical Spin$^c$ structure with $L_{\det} = \mathcal{O}(3)$.

%---------------------------------------------------------------------
\subsection{Step 4: Twist bundle and $k_{\max}$ (Bridge Lemma)}
%---------------------------------------------------------------------

The twist bundle over $\mathbb{C}P^2$ is chosen as:
\begin{equation}\tag{E4a}
E = \mathcal{O}(9) \oplus \mathcal{O}^{\oplus 5}.
\end{equation}

The selection $a = 9$ arises from $a = q_1^2 \cdot q_1 = 3^2 \cdot 1 = 9$
where $q_1 = 3$ is the minimal U(1) flux quantum from anomaly
cancellation/Spin$^c$ constraint (Lemma F.5).  The $n = 5$ trivial
summands come from the $(3,2,1)$ partition structure ($N = 3$ colors,
$n = N + 2 = 5$ residual channels).

The maximum Chern--Simons level is the closed Spin$^c$ index:
\begin{equation}\tag{E4b}
k_{\max} = \chi(\mathbb{C}P^2, E) = \chi(\mathcal{O}(9)) + 5\,\chi(\mathcal{O}).
\end{equation}

Using the Hirzebruch--Riemann--Roch formula on $\mathbb{C}P^2$:
\begin{equation}\tag{E4c}
\chi(\mathbb{C}P^2, \mathcal{O}(m)) = \binom{m+2}{2} \quad \text{for } m \geq 0.
\end{equation}

Therefore:
\begin{align}
\chi(\mathcal{O}(9)) &= \binom{11}{2} = 55, \tag{E4d}\\
\chi(\mathcal{O})     &= \binom{2}{2} = 1, \tag{E4e}
\end{align}
and:
\begin{equation}\tag{E4f}
\boxed{k_{\max} = 55 + 5 = 60.}
\end{equation}

\textit{Status:} Rigorous (Hirzebruch--Riemann--Roch).

\textit{Consistency checks:}
\begin{itemize}[nosep]
\item Icosahedral: $k_{\max} = 60 = |A_5|$ (McKay correspondence).
\item $E_8$ echo: $\mathrm{roots}(E_8)/4 = 240/4 = 60$.
\item Finite-symmetry closure: $A_5$ is the minimal simple group acting
  faithfully on $\mathbb{R}^3$ with exactly 5 conjugacy classes.
\end{itemize}

%---------------------------------------------------------------------
\subsection{Step 5: Chern--Simons weights on $S^3$}
%---------------------------------------------------------------------

The $S^3$ factor supports Chern--Simons gauge theory.  For U(1) gauge
fields, the Chern--Simons action is:
\begin{equation}\tag{E5a}
S_{\mathrm{CS}} = \frac{k}{4\pi}\int_{S^3} A \wedge dA.
\end{equation}

The SU(2)$_k$ Chern--Simons partition function on $S^3$ gives the weight:
\begin{equation}\tag{E5b}
w(k) = \frac{2}{k+2}\sin^2\!\frac{\pi}{k+2}, \quad k = 0, 1, \ldots, k_{\max} - 1.
\end{equation}

\textit{Status:} Derived from the SU(2) WZW/CS partition function
$Z_{SU(2)_k}(S^3) = \sqrt{2/(k+2)}\,\sin(\pi/(k+2))$, with
$w(k) = |Z|^2$.

%---------------------------------------------------------------------
\subsection{Step 6: The CS weighted average $\beta_{U(1)}$}
%---------------------------------------------------------------------

The effective U(1) lattice coupling is defined as the weighted expectation:
\begin{equation}\tag{E6a}
\beta_{U(1)} = \langle k + 2 \rangle
= \frac{\displaystyle\sum_{k=0}^{k_{\max}-1}(k+2)\,w(k)}
       {\displaystyle\sum_{k=0}^{k_{\max}-1}w(k)}.
\end{equation}

\textbf{Numerical evaluation at $k_{\max} = 60$:}
\begin{equation}\tag{E6b}
\beta_{U(1)}\big|_{k_{\max}=60} = \frac{8.3245\ldots}{2.1924\ldots} = 3.796945\ldots \approx 3.80.
\end{equation}

\textit{Status:} Computable pure number.  Verified by Python/Mathematica.

%---------------------------------------------------------------------
\subsection{Step 7: Heat kernel on $S^3$ (spectral data)}
%---------------------------------------------------------------------

The heat kernel on $S^3$ with radius $R$ has the spectral expansion:
\begin{equation}\tag{E7}
K(t; S^3) = \sum_{n=0}^{\infty}(n+1)^2\,e^{-n(n+2)t/R^2}.
\end{equation}

The factor $(n+1)^2$ is the degeneracy of the $n$-th eigenvalue
$\lambda_n = n(n+2)/R^2$ of the scalar Laplacian on $S^3$.

\textit{Status:} Standard spectral geometry result.

%---------------------------------------------------------------------
\subsection{Step 8: Stiffness ratio from Ricci curvature}
%---------------------------------------------------------------------

From the Frame Stiffness Theorem (Theorem F.13 in Appendix F), the
ratio of U(1) to SU(2) stiffness coefficients is:
\begin{equation}\tag{E8}
\frac{\kappa_{U(1)}}{\kappa_{SU(2)}} = \frac{n_1}{n_2} = \frac{1}{2}.
\end{equation}

\textit{Status:} Derived from the Ricci curvature structure of the
internal manifold.  Lattice-verified: measured ratio $0.495 \pm 0.020$.

%---------------------------------------------------------------------
\subsection{Step 9: Wilson ratio from topology}
%---------------------------------------------------------------------

The Wilson action ratio is derived (not a convention):
\begin{equation}\tag{E9a}
\frac{\beta_{SU(2)}}{\beta_{U(1)}} = \frac{n_2}{n_1} \times N_{\mathrm{gen}}
= 2 \times 3 = 6.
\end{equation}

Therefore:
\begin{equation}\tag{E9b}
\beta_{SU(2)} = 6 \times 3.7969 = 22.782 \approx 22.80.
\end{equation}

\textit{Status:} Derived.  Lattice-verified: 10 Wilson ratios tested
(3--9 including fractional), only exactly 6 yields $\alpha = 1/137$.

%---------------------------------------------------------------------
\subsection{Step 10: Electroweak mixing formula}
%---------------------------------------------------------------------

The gauge couplings are extracted from stiffness coefficients via
Wilson normalization:
\begin{align}
g_1^2 &= \frac{1}{\kappa_{U(1)}}, \tag{E10a}\\
g_2^2 &= \frac{4}{\kappa_{SU(2)}}. \tag{E10b}
\end{align}

The electromagnetic coupling satisfies:
\begin{equation}\tag{E10c}
\frac{1}{e^2} = \frac{1}{g_1^2} + \frac{1}{g_2^2}
= \kappa_{U(1)} + \frac{\kappa_{SU(2)}}{4}.
\end{equation}

Using the stiffness ratio $\kappa_{SU(2)} = 2\kappa_{U(1)}$ (from E8):
\begin{equation}\tag{E10d}
\frac{1}{e^2} = \kappa_{U(1)} + \frac{2\kappa_{U(1)}}{4}
= \frac{3}{2}\,\kappa_{U(1)}.
\end{equation}

Since $\alpha = e^2/(4\pi)$:
\begin{equation}\tag{E10e}
\boxed{\alpha^{-1} = \frac{4\pi}{e^2} = 4\pi \cdot \frac{3}{2}\,\kappa_{U(1)}
= 6\pi\,\kappa_{U(1)}.}
\end{equation}

This is the \textbf{master formula}.  The entire derivation reduces to
predicting $\kappa_{U(1)}$.

\textit{Status:} Exact, from standard electroweak mixing with the DFD
stiffness ratio.  Requires $\kappa_{U(1)} = 7.26999\ldots$ for the
experimental value.

%---------------------------------------------------------------------
\subsection{Step 11: Lattice Monte Carlo measurement (Route 1 completion)}
%---------------------------------------------------------------------

At the derived parameter point $(\beta_{U(1)}, \beta_{SU(2)}) = (3.80, 22.80)$,
lattice Monte Carlo simulations measure the renormalized stiffnesses via the
background-field method.

\textbf{Key result} (86 lattice runs, $L = 4$--$16$):
\begin{equation}\tag{E11a}
\alpha_W = 0.007297 \pm 0.000094.
\end{equation}

Therefore:
\begin{equation}\tag{E11b}
\alpha^{-1} = 1/0.007297 = 137.04 \pm 1.8 \quad (\sim 1\%\text{ precision}).
\end{equation}

\textit{Status:} Verified by independent lattice runs across $L = 6$--$16$.
L12 best result: $\alpha_W = 0.007291$ ($-0.08\%$ from experiment).

\begin{tcolorbox}[colback=green!5!white,colframe=green!50!black,
title={Route 1 Result}]
$\alpha^{-1} \approx 137.0 \pm 1.8$ from first-principles inputs +
lattice MC, with zero continuous free parameters.
\end{tcolorbox}

%=============================================================================
\section{Route 2: Convention-Locked Spectral Action $\to$ Sub-ppm $\alpha$}
\label{sec:route2}
%=============================================================================

Route 2 bypasses the lattice and computes $\kappa_{U(1)}$ analytically
from the spectral action on the Toeplitz-truncated internal geometry.
It uses Steps 1--10 above plus five additional locked conventions
(Steps 12--16).

%---------------------------------------------------------------------
\subsection{Step 12: Operator choice (locked)}
%---------------------------------------------------------------------

On the internal microsector $X = \mathbb{C}P^2 \times S^3$, the operator
is the connection Laplacian:
\begin{equation}\tag{E12}
P = -g^{ij}\nabla_i\nabla_j,
\end{equation}
acting on the internal bundle carrying emergent gauge degrees of freedom.

The U(1) factor is implemented via twisting by a line bundle over
$\mathbb{C}P^2$ with curvature proportional to the K\"ahler form $\omega$.
Since $\nabla\omega = 0$ (the K\"ahler form is parallel), derivative terms
in higher Seeley--DeWitt coefficients vanish automatically.

\textit{Status:} Locked.  Alternative operators introduce ppm-level
ambiguities; this is the unique choice without hidden knobs.

%---------------------------------------------------------------------
\subsection{Step 13: Spectral action with plateau cutoff (locked)}
%---------------------------------------------------------------------

The spectral action is:
\begin{equation}\tag{E13a}
S = \mathrm{Tr}\,f(P/\Lambda^2),
\end{equation}
where $f$ satisfies the \textbf{plateau condition}: $f$ is constant
near the origin, so $f^{(n)}(0) = 0$ for all $n \geq 1$:
\begin{equation}\tag{E13b}
f_{-2} = f_{-4} = \cdots = 0.
\end{equation}

This eliminates all subleading heat-kernel contributions ($a_6, a_8, \ldots$),
leaving only the $a_4$ gauge kinetic term.

\textit{Status:} Locked.  This is the unique cutoff that prevents
$a_6$ and higher from acting as hidden tuning knobs (with generic smooth
cutoffs, the $a_6$ contribution would be $\sim 2\%$).

%---------------------------------------------------------------------
\subsection{Step 14: Toeplitz truncation and the $(k+3)/(k+4)$ factor (locked)}
%---------------------------------------------------------------------

The finite-$k$ truncation is implemented via Toeplitz quantization at
level $m = k + 3$ on $\mathbb{C}P^1 \subset \mathbb{C}P^2$.

The shift $m = k + 3$ arises from the Spin$^c$ structure:
\begin{equation}\tag{E14a}
K_{\mathbb{C}P^2} = \mathcal{O}(-3) \quad\Rightarrow\quad
L_{\det} = K^{-1} = \mathcal{O}(3).
\end{equation}

When restricting to $\mathbb{C}P^1 \subset \mathbb{C}P^2$, the line
bundle $\mathcal{O}(k) \otimes L_{\det}$ becomes $\mathcal{O}(k+3)$,
giving:
\begin{equation}\tag{E14b}
d = \dim H^0(\mathbb{C}P^1, \mathcal{O}(m)) = m + 1 = k + 4.
\end{equation}

For $k = k_{\max} = 60$:
\begin{equation}\tag{E14c}
d = 60 + 4 = 64.
\end{equation}

The determinant-channel removal at finite $d$ fixes the spectral cutoff:
\begin{equation}\tag{E14d}
\Lambda^3 = k \cdot \frac{d-1}{d} = k \cdot \frac{k+3}{k+4}.
\end{equation}

For $k = 60$:
\begin{equation}\tag{E14e}
\Lambda^3 = 60 \times \frac{63}{64} = 60 \times 0.984375 = 59.0625.
\end{equation}

\textit{Note:} The text at line 210 of \texttt{section\_alpha\_locked.tex}
also refers to a ``baseline normalization $\Lambda^3 = 885.9375$ (from
$k = 60$, $a = 9$, $n = 5$, $N = 3$).''  This is $885.9375 = 60 \times (63/64) \times 15
= 59.0625 \times 15$, where the factor 15 encodes the product
$a \cdot n / N = 9 \times 5 / 3 = 15$ from the twist bundle data.
The relationship between $\Lambda^3 = 59.0625$ and $\Lambda^3 = 885.9375$
depends on whether the twist-bundle multiplicity is folded into $\Lambda^3$
or kept as a separate prefactor.

\textit{Status:} Locked.  The $(k+3)/(k+4) = 63/64$ factor is uniquely
fixed by the truncation rule.

%---------------------------------------------------------------------
\subsection{Step 15: SM content inputs}
%---------------------------------------------------------------------

Three Standard Model data enter through the $a_4$ trace over the
internal fermion Hilbert space:

\begin{equation}\tag{E15a}
N_{\mathrm{species}} = 7 \quad \text{(SM SU(2) doublet components)}.
\end{equation}

\begin{equation}\tag{E15b}
\mathrm{Tr}(Y^2) = 10 \quad \text{(SM hypercharge sum over one generation)}.
\end{equation}

\begin{equation}\tag{E15c}
g_F = 8 \quad \text{(spectral triple grading: } J \times \gamma \times \mathbb{C}\text{)}.
\end{equation}

These combine into the hypercharge weighting:
\begin{equation}\tag{E15d}
w = \frac{N_{\mathrm{species}}}{g_F \cdot \mathrm{Tr}(Y^2)}
= \frac{7}{8 \times 10} = \frac{7}{80} = 0.0875.
\end{equation}

\textit{Status:} $N_{\mathrm{species}}$ and $\mathrm{Tr}(Y^2)$ are SM
content (empirical input).  $g_F = 8$ is derived from the spectral triple
structure.

%---------------------------------------------------------------------
\subsection{Step 16: The forced microsector fork (BOOST factor)}
%---------------------------------------------------------------------

Two microsector trace choices exist:

\textbf{Branch A (Regular-Module --- survives):}
\begin{equation}\tag{E16a}
\mathcal{H}_F = A = M_d(\mathbb{C}), \quad
\dim(\mathcal{H}_F) = d^2 = 4096.
\end{equation}

UV-normalized trace:
\begin{equation}\tag{E16b}
\mathrm{tr}_{\mathrm{dem}}(\cdot) = \frac{1}{d^2}\,\mathrm{Tr}_{\mathcal{H}_F}(\cdot).
\end{equation}

Physics-normalized trace on $\mathfrak{su}(d)$:
\begin{equation}\tag{E16c}
\mathrm{tr}_{\mathfrak{su}}(\cdot) = \frac{1}{d^2 - 1}\,\mathrm{Tr}_{\mathfrak{su}}(\cdot).
\end{equation}

The forced conversion factor is:
\begin{equation}\tag{E16d}
\boxed{\varepsilon_{\mathrm{adj}}^{(A)} = \frac{d^2}{d^2 - 1}
= \frac{4096}{4095} = 1.000244\ldots}
\end{equation}

\textbf{Branch B (Fermion-Rep --- falsified):}
\begin{equation}\tag{E16e}
\varepsilon_{\mathrm{adj}}^{(B)} = \frac{d^2 - 1}{d^2}
= \frac{4095}{4096} = 0.999756\ldots
\end{equation}

\textit{Status:} The regular-module choice is forced by the no-knobs
policy (Branch B misses by 43 ppm, unrescuable).

%---------------------------------------------------------------------
\subsection{Step 17: Assembly into $\kappa_{U(1)}$ and $\alpha^{-1}$}
%---------------------------------------------------------------------

The spectral action on the Toeplitz-truncated $\mathbb{C}P^2 \times S^3$
with all locked components gives the stiffness coefficient:

\begin{equation}\tag{E17a}
\kappa_{U(1)}^{\mathrm{pred}} = 7.26999\ldots
\end{equation}

Therefore, by the master formula (E10e):
\begin{equation}\tag{E17b}
\boxed{\alpha^{-1} = 6\pi \times 7.26999\ldots = 137.03599985.}
\end{equation}

Residual:
\begin{equation}\tag{E17c}
\frac{\alpha^{-1}_{\mathrm{DFD}} - \alpha^{-1}_{\mathrm{exp}}}
     {\alpha^{-1}_{\mathrm{exp}}}
= \frac{137.03599985 - 137.035999084}{137.035999084}
= -0.006\text{ ppm}.
\end{equation}

\begin{tcolorbox}[colback=green!5!white,colframe=green!50!black,
title={Route 2 Result}]
$\alpha^{-1} = 137.03599985$ (residual $-0.006$ ppm from experiment).

All conventions locked. Zero continuous free parameters.
\end{tcolorbox}

%=============================================================================
\section{Complete Equation Inventory}
\label{sec:inventory}
%=============================================================================

\renewcommand{\arraystretch}{1.25}
\begin{table}[h]
\centering
\caption{Every equation in the derivation chain, with source file and line.}
\label{tab:inventory}
\resizebox{\textwidth}{!}{%
\begin{tabular}{clllll}
\toprule
\textbf{Tag} & \textbf{Equation} & \textbf{Value at $k=60$} & \textbf{Source file} & \textbf{Line} & \textbf{Status} \\
\midrule
E1 & $\mathcal{M}_7 = \mathbb{C}P^2 \times S^3$ & --- & appendix\_K.tex & 12 & Postulate \\
E2 & $K_{\mathbb{C}P^2} = \mathcal{O}(-3)$ & $-3$ & section\_alpha\_locked.tex & 72 & Rigorous \\
E3 & $L_{\det} = \mathcal{O}(3)$ & $+3$ & section\_alpha\_locked.tex & 72 & Rigorous \\
E4a & $E = \mathcal{O}(9) \oplus \mathcal{O}^{\oplus 5}$ & --- & appendix\_K.tex & 59 & Derived ($q_1=3$) \\
E4c & $\chi(\mathcal{O}(m)) = \binom{m+2}{2}$ & --- & appendix\_K.tex & 62 & Rigorous (HRR) \\
E4f & $k_{\max} = 55 + 5 = 60$ & 60 & appendix\_K.tex & 68 & Rigorous \\
E5b & $w(k) = \frac{2}{k+2}\sin^2\!\frac{\pi}{k+2}$ & --- & appendix\_K.tex & 35, 74 & Derived (CS) \\
E6a & $\beta_{U(1)} = \langle k+2 \rangle$ & 3.7969 & appendix\_K.tex & 33, 78 & Computed \\
E7 & $K(t;S^3) = \sum (n+1)^2 e^{-n(n+2)t/R^2}$ & --- & appendix\_K.tex & 41 & Standard \\
E8 & $\kappa_{U(1)}/\kappa_{SU(2)} = 1/2$ & 0.5 & appendix\_K.tex & 127 & Derived (Thm F.13) \\
E9a & Wilson ratio $= 6$ & 6 & appendix\_K.tex & 133 & Derived \\
E10e & $\alpha^{-1} = 6\pi\,\kappa_{U(1)}$ & --- & Alpha\_Inverse\_Formula.tex & 134 & Exact \\
E12 & $P = -g^{ij}\nabla_i\nabla_j$ & --- & section\_alpha\_locked.tex & 26 & Locked \\
E13a & $S = \mathrm{Tr}\,f(P/\Lambda^2)$ & --- & section\_alpha\_locked.tex & 42 & Locked \\
E13b & $f_{-2} = f_{-4} = \cdots = 0$ & 0 & section\_alpha\_locked.tex & 49 & Locked \\
E14b & $d = k + 4$ & 64 & section\_alpha\_locked.tex & 67 & Rigorous \\
E14d & $\Lambda^3 = k(k+3)/(k+4)$ & 59.0625 & section\_alpha\_locked.tex & 79 & Derived \\
E15d & $w = 7/80$ & 0.0875 & section\_alpha\_locked.tex & 200 & SM + Derived \\
E16d & $\varepsilon_{\mathrm{adj}} = 4096/4095$ & 1.000244 & section\_alpha\_locked.tex & 114 & Forced \\
E17a & $\kappa_{U(1)} = 7.26999\ldots$ & 7.26999 & (numerical) & --- & \textbf{See Gap 1} \\
E17b & $\alpha^{-1} = 137.03599985$ & 137.036 & section\_alpha\_locked.tex & 145, 253 & Result \\
\bottomrule
\end{tabular}%
}
\end{table}

%=============================================================================
\section{Identified Gaps in the Published Chain}
\label{sec:gaps}
%=============================================================================

\begin{tcolorbox}[colback=red!5!white,colframe=red!60!black,
title={GAP 1: The $a_4$ computation is never shown explicitly}]
\textbf{What is missing:} The derivation claims $\kappa_{U(1)} = 7.26999\ldots$
from the $a_4$ Seeley--DeWitt coefficient of the connection Laplacian on
Toeplitz-truncated $\mathbb{C}P^2 \times S^3$.  But the standard
Seeley--DeWitt--Gilkey formula
\begin{equation*}
a_4(P) = \frac{1}{(4\pi)^{d/2}} \frac{1}{360} \int_X \mathrm{tr}_V\!\Big(
60ER + 180E^2 + 30\,\Omega_{\mu\nu}\Omega^{\mu\nu}
+ (5R^2 - 2R_{\mu\nu}R^{\mu\nu} + 2R_{\mu\nu\rho\sigma}R^{\mu\nu\rho\sigma})\mathrm{Id}_V
\Big)\mathrm{dvol}
\end{equation*}
is \textbf{never written out, evaluated, or cited with its value} on the
specific geometry.  The gauge kinetic term emerges from the
$\Omega_{\mu\nu}\Omega^{\mu\nu}$ piece, but this evaluation is absent.

\textbf{Inputs to this computation:}
\begin{itemize}[nosep]
\item $\mathbb{C}P^2$: K\"ahler--Einstein, $R = 24$, $R_{\mu\nu} = 6g_{\mu\nu}$,
  $R_{\mu\nu\rho\sigma}R^{\mu\nu\rho\sigma} = 24$ (unit FS metric).
\item $S^3$: constant curvature, $R = 6/r^2$,
  $R_{\mu\nu}R^{\mu\nu} = R_{\mu\nu\rho\sigma}R^{\mu\nu\rho\sigma} = 12/r^4$.
\item Bundle curvature: $\Omega = i\omega$ on $\mathbb{C}P^2$, zero on $S^3$.
\item Product formula: $a_4(X) = \sum_{p+q=4} a_p(\mathbb{C}P^2) \cdot a_q(S^3)$.
\item Toeplitz modification: replaces continuous integral with finite matrix
  trace, introducing the $(63/64)$ and $(4096/4095)$ factors.
\end{itemize}

\textbf{Output:} $\kappa_{U(1)} = 7.26999\ldots$

The reader is asked to accept that these eight inputs, when fed through the
$a_4$ heat-kernel machinery with Toeplitz truncation, produce the stated
value of $\kappa_{U(1)}$.  The explicit computation connecting inputs to
output has not been published.
\end{tcolorbox}

\begin{tcolorbox}[colback=orange!5!white,colframe=orange!60!black,
title={GAP 2: How the eight Table XXVIII components combine}]
\textbf{What is missing:} An explicit formula of the form
\begin{equation*}
\kappa_{U(1)} = F\!\left(\Lambda^3,\; w,\; \varepsilon_{\mathrm{adj}},\;
\text{curvature data}\right) = 7.26999\ldots
\end{equation*}
does not appear in the DFD corpus.  The text states ``All above combined''
in Table XXVIII but never writes the combining formula.

\textbf{Best available approximation}
(\texttt{Alpha\_Inverse\_Formula.tex}, Sec.~4):
\begin{equation*}
\alpha^{-1} \approx \frac{35\pi}{3}\;\beta_{U(1)}\;\frac{d-1}{d}
\;\frac{d^2}{d^2-1}
\end{equation*}
which gives $137.024$ (85 ppm residual).  The exact sub-ppm formula
requires the full spectral action computation (Gap 1).
\end{tcolorbox}

\begin{tcolorbox}[colback=orange!5!white,colframe=orange!60!black,
title={GAP 3: Toeplitz modification of $a_4$ not derived from first principles}]
How the Toeplitz truncation at level $d = 64$ modifies the continuous
$a_4$ integral is \textit{stated} (the $(k+3)/(k+4)$ factor and boost)
but \textit{not derived} from Toeplitz symbol calculus.  The assertion
is that this is the ``unique finite-size factor permitted by the
truncation rule,'' but the derivation of this uniqueness is absent.
\end{tcolorbox}

\begin{tcolorbox}[colback=orange!5!white,colframe=orange!60!black,
title={GAP 4: $C_{\mathrm{loop}} = 1/8$ status}]
The relation $k_a = 3/(8\alpha)$ (Theorem F.X) uses a loop factor
$C_{\mathrm{loop}} = 1/8$ that is described in
\texttt{section\_alpha\_relations.tex} as ``plausible from heat kernel
structure but requires explicit verification'' (status: Plausible B).
This factor should emerge from the $a_4$ computation of Gap 1.
\end{tcolorbox}

%=============================================================================
\section{What Is NOT a Gap}
\label{sec:not-gaps}
%=============================================================================

The following are sometimes suspected as gaps but are fully present
in the published sources:

\begin{enumerate}
\item \textbf{$k_{\max} = 60$:} Rigorously derived from HRR
  (E4c--E4f).  Also confirmed by lattice scanning (Appendix K,
  Table~\ref{tab:kmax-discovery} of source) and icosahedral closure.

\item \textbf{$\beta_{U(1)} = 3.7969$:} Explicitly computed as a
  definite sum (E6a--E6b).  Python-verifiable.

\item \textbf{Wilson ratio $= 6$:} Derived from topology (E9a).
  Lattice-verified by 10-ratio scan.

\item \textbf{Stiffness ratio $= 1/2$:} Derived (Theorem F.13).
  Lattice-measured: $0.495 \pm 0.020$.

\item \textbf{$\varepsilon_{\mathrm{adj}} = 4096/4095$:} Explicitly
  derived from the regular-module trace structure (E16a--E16d).

\item \textbf{Master formula $\alpha^{-1} = 6\pi\kappa_{U(1)}$:}
  Derived from standard electroweak mixing (E10a--E10e).

\item \textbf{Lattice verification:} 86 runs, multiple lattice sizes,
  systematic independence checks ($k_0$, $\varepsilon$, Wilson ratio
  scan, $\beta$ bracket).
\end{enumerate}

%=============================================================================
\section{Summary: The Derivation Chain at a Glance}
\label{sec:summary}
%=============================================================================

\begin{center}
\fbox{\parbox{0.92\textwidth}{
\textbf{Layer 0: Geometry}

$\mathbb{C}P^2 \times S^3$ $\xrightarrow{K = \mathcal{O}(-3)}$
Spin$^c$, $L_{\det} = \mathcal{O}(3)$
$\xrightarrow{q_1 = 3}$ $E = \mathcal{O}(9) \oplus \mathcal{O}^{\oplus 5}$

$\downarrow$ HRR

$k_{\max} = \chi(E) = 55 + 5 = 60$

\medskip
\textbf{Layer 1: Chern--Simons}

$w(k) = \frac{2}{k+2}\sin^2\!\frac{\pi}{k+2}$
$\xrightarrow{\text{weighted avg}}$
$\beta_{U(1)} = 3.7969$

Wilson ratio $= 6$ $\to$ $\beta_{SU(2)} = 22.80$

\medskip
\textbf{Layer 2: Stiffness and Electroweak Mixing}

$\kappa_{U(1)}/\kappa_{SU(2)} = 1/2$ (Frame Stiffness Thm)

$\alpha^{-1} = 6\pi\,\kappa_{U(1)}$ (Master Formula)

\medskip
\textbf{Layer 3a (Route 1): Lattice MC}

MC at $(3.80, 22.80)$ $\to$ $\kappa_{U(1)} \approx 7.27$
$\to$ $\alpha^{-1} \approx 137$ ($\pm 1\%$)

\medskip
\textbf{Layer 3b (Route 2): Spectral Action}

Connection Laplacian + Plateau cutoff + Toeplitz at $d=64$

+ SM content ($N_{\mathrm{sp}}=7$, $\mathrm{Tr}(Y^2)=10$, $g_F=8$, $w=7/80$)

+ BOOST $4096/4095$

$\to$ $\kappa_{U(1)} = 7.26999\ldots$
$\to$ $\alpha^{-1} = 137.03599985$ (sub-ppm)

\medskip
\textcolor{red}{\textbf{GAP:}} The explicit computation connecting
``$a_4$ on Toeplitz-truncated geometry + 8 inputs'' $\to$
``$\kappa_{U(1)} = 7.26999$'' is asserted but never shown.
}}
\end{center}

%=============================================================================
\section{Numerical Cross-Check: Approximate Closed Form}
\label{sec:approx}
%=============================================================================

The best available closed-form approximation is:
\begin{equation}\tag{APPROX}
\alpha^{-1} \approx \frac{35\pi}{3}\;\beta_{U(1)}\;\frac{d-1}{d}
\;\frac{d^2}{d^2-1}
\end{equation}

where the prefactor decomposes as:
\begin{equation}
\frac{35\pi}{3} = 4\pi \times \underbrace{\frac{5}{2}}_{\text{EW mixing}}
\times \underbrace{\frac{7}{6}}_{\text{spectral triple}}.
\end{equation}

\textbf{Numerical evaluation:}
\begin{align}
\frac{35\pi}{3} &= 36.652\ldots \\
\beta_{U(1)} &= 3.7969 \\
\frac{d-1}{d} &= \frac{63}{64} = 0.984375 \\
\frac{d^2}{d^2-1} &= \frac{4096}{4095} = 1.000244\ldots \\
\text{Product} &= 36.652 \times 3.7969 \times 0.984375 \times 1.000244 \\
&= 137.024 \quad (\text{residual: } 85\text{ ppm from experiment}).
\end{align}

The 85 ppm residual of this approximation vs.\ the exact sub-ppm
result reflects non-perturbative corrections in the spectral action
computation that the simple algebraic product does not capture.

%=============================================================================
\section{Comparison Table: Both Branches}
\label{sec:branches}
%=============================================================================

\begin{table}[h]
\centering
\caption{Microsector fork: numerical comparison at $k = 60$, $d = 64$.}
\renewcommand{\arraystretch}{1.3}
\begin{tabular}{lccc}
\toprule
Branch & Factor & $\alpha^{-1}$ & Residual (ppm) \\
\midrule
\textbf{A (regular-module)} & $4096/4095$ & $137.03599985$ & $-0.006$ \\
B (fermion-rep) & $4095/4096$ & $137.03014445$ & $+42.7$ \\
\midrule
Experimental & --- & $137.035999084(21)$ & --- \\
\bottomrule
\end{tabular}
\end{table}

Branch B cannot be rescued: $\Delta w/w = -200\%$, $\Delta g_F/g_F = +200\%$,
or $a_6$ tuning with $f_{-2}/f_0 \sim 10^{-3}$ would all be required,
violating the no-knobs policy.

\end{document}
